\begin{papercard}{harkous2018}
\cf{Stages} \stagetag{L} (automatic annotation of policies, evaluated against expert labels), \stagetag{I} (human--machine agreement and two-author agreement reported with accuracy and Cohen's $\kappa$), \stagetag{C} (re-use of an existing taxonomy), \stagetag{D} (corpus and test-set construction). Polisis is a deep-learning framework that annotates privacy-policy segments with the OPP-115 categories and attribute values, and serves two applications: automatic privacy icons and PriBot, a free-form question-answering system. This is the pre-LLM reference for ``how well can a machine do what annotators do on policy text'', and it shows how such systems are evaluated.

\cf{Problem and aim} Policies are long (McDonald \& Cranor's estimate: 201 hours per year to read all those encountered) and short notices (nutrition labels, icons) need human effort that does not scale \pc{harkous2018}{2}{In 2008, McDonald and Cranor [4] estimated that it would take an average user 201 hours to read all the privacy policies encountered in a year}. Polisis predicts for each segment ``a set of classes that account for both the high-level aspects and the fine-grained classes'' and is explicitly ``not intended to replace the privacy policy as a legal document'' \pc{harkous2018}{3}{Polisis is not intended to replace the privacy policy – as a legal document – with an automated interpretation}. Claims: 88.4\,\% accuracy on automatic icon assignment, PriBot has a correct answer in its top three for 82\,\% of test questions (68\,\% at rank 1) and 89\,\% of top-3 sets contain an answer users find relevant.

\cf{Architecture} Three layers: application layer (queries), data layer (segmenter) and ML layer. A policy is a list of ``semantically coherent segments'' \pc{harkous2018}{3}{Polisis treats a privacy policy as a list of semantically coherent segments}. \emph{Segmenter}: headless Chrome loads the page; irrelevant HTML (scripts, headers, menus) is removed; the scraped text covers on average 99.08\,\% of a manually fetched version for the top 200 Alexa sites \pc{harkous2018}{4}{the segmenter’s scraped policy covers, on average, 99.08\% of the content of the manually fetched policy}; list items of up to 20 words are merged with the list's introductory sentence, longer ones become paragraphs prefixed by that sentence \pc{harkous2018}{4}{For short list items (maximum of 20 words per element)}; coarse splitting on div/p tags, then GraphSeg (embedding-based semantic cliques) for long pieces.

\cf{Machine-learning layer} (1) \emph{Domain-specific embeddings}: fastText trained on a corpus of 130,326 policies obtained from links of 199,186 policy URLs among 1.4 million Google Play apps (PlayDrone); subword vectors handle unknown and misspelled words \pc{harkous2018}{5}{We crawled the metadata of more than 1.4 million Android apps available via the PlayDrone project [27] to find the links to 199,186 privacy policies}. (2) \emph{Training data}: OPP-115 (115 policies, 23K annotated practices; 20 attributes and 138 values, 122 of them with more than 20 labels) \pc{harkous2018}{5}{In total, there were 20 distinct attributes and 138 distinct values across all attributes}. (3) \emph{Hierarchical multi-label classification}: a category-level classifier (10 categories, present if probability $\ge0.5$) and attribute-level classifiers for each attribute's values; attribute outputs only count under the dominant category; CNN with frozen embeddings, ReLU, max-pooling, two dense layers, sigmoid and multi-label cross-entropy; 20 attribute classifiers plus two category classifiers; 65 policies for training and 50 for testing; hyper-parameters by randomised grid search \pc{harkous2018}{6}{For training the classifiers, we used the data from 65 policies in the OPP-115 dataset, and we kept 50 policies as a testing set}. Metrics are macro-averaged over predicting presence \emph{and} absence of each label, since a model predicting everything present would otherwise score perfectly \pc{harkous2018}{6}{A successful multilabel classifier should not only predict the presence of a label, but also its absence}. Category-level results for user queries (Table 1): averages precision .87, recall .83, F1 .84; best Do Not Track (.97) and Specific Audiences (.95), weakest Data Retention (F1 .71) and User Choice/Control (.75); the authors state these are significantly higher than the baselines of the OPP-115 paper \pc{harkous2018}{7}{these metrics are significantly higher than the models presented in the original OPP-115 paper}.

\cf{Application 1: privacy icons} Structured queries are first-order-logic predicates over predicted classes (e.g.\ information type = location $\wedge$ purpose = marketing $\wedge$ choice = opt-out) \pc{harkous2018}{7}{A structured query is a combination of first-order logic predicates over the predicted privacy classes and the policy segments}. The Disconnect icons (Expected Use, Expected Collection, Precise Location, Data Retention, Children Privacy) have imprecise rules, so the authors translated them into explicit label rules (Table 2) as ``best effort'' \pc{harkous2018}{8}{Given that the rules behind the Disconnect icons are not precisely defined, we translated their description into explicit first-order logic queries}. \emph{Evaluation}: on the 50 held-out policies, icons computed from Polisis labels are compared with icons computed from the \emph{union} of the three law-student annotations; accuracy per icon 92\,\%, 88\,\%, 84\,\%, 80\,\%, 98\,\% (mean \textbf{88.4\,\%}) and Cohen's $\kappa$ .76, .69, .68, .63, .95, i.e.\ ``substantial to almost perfect'' on the Landis \& Koch scale; Hellinger distances between colour distributions .02--.21 \pc{harkous2018}{9}{The average accuracy across icons is 88.4\%}. For context they cite human judges assessing policies at 88.3--98.3\,\% (mean 92.7\,\%) agreement.

\cf{Auditing a certifier} 1,010 policies from Dec 2013--Jan 2014 (Ramanath et al.), 354 with Disconnect icons; with the conservative rule the Disconnect colour distribution differs sharply from Polisis' (Hellinger .71 and .61, 3--5 times the evaluation distance); only after assigning yellow to any policy with an opt-out segment anywhere (``very permissive'') do the distributions come near (.22 and .19). Interpretation: the TRUSTe-derived Disconnect icons may be lenient, though the authors admit possible version mismatch between policies \pc{harkous2018}{10}{Only in this interpretation of the icons would the distributions of Disconnect and Polisis come within reasonable proximity}. Human oversight stays in the loop where legal liability could arise \pc{harkous2018}{10}{by automatically generating icons, we do not intend to push humans completely out of the loop}.

\cf{Application 2: PriBot} Question answering is cast as ranking policy segments by a proximity score computed from the predicted category/value probabilities of question and segment (category probability squared, normalised dot product, weighted by an entropy-based certainty), plus a confidence indicator combining score, certainty and fraction of known words; a ``does vs.\ does not'' classifier (cross-validated F1 95\,\%) highlights conflicting answers \pc{harkous2018}{11}{we used the same CNN classifier of Sec. 4 and exploited the fact that the OPP-115 dataset had optional labels}. \emph{Test set}: questions that Twitter users tweeted at companies. Queries found reply tweets pointing to a policy (``our privacy policy''), back-tracked to the parent question: 4,743 pairs \pc{harkous2018}{12}{a set of 4,743 pairs of tweets}; automatic filter (question mark, at least four content words, a policy link) left 260; two authors independently labelled each valid/invalid, agreeing on 231 (182 valid, 49 invalid), $\kappa=.84$ \pc{harkous2018}{12}{The two annotators agreed on 231 of the question tweets (of the 260), tagging 182 as valid and 49 as invalid}; 120 valid pairs sampled. \emph{Ground truth}: the two authors independently selected best-matching segments (mean 45 segments per policy) before seeing model output, then consolidated; they agreed on at least one answer in 88\,\% of the questions with matches, $\kappa=.65$ \pc{harkous2018}{13}{The annotators agreed on at least one answer in 88\% of the questions for which they found matching segments}. \emph{Baselines}: Retrieval (BM25 with idf from the 130K corpus), SemVec (one 81-class classifier's semantic vectors, Euclidean distance), Random. \emph{Results}: top-1 PriBot .68 vs.\ Retrieval .39, SemVec .27, Random .08 (Fisher's exact tests with Bonferroni); top-2/3/4 .75/.82/.87 \pc{harkous2018}{13}{PriBot has a top-k score of 0.68, which is significantly larger than the scores of 0.39 (Retrieval), 0.27 (SemVec), and 0.08 (Random)}; NDCG and MAP by policy length; wrong answers receive lower confidence (mean .37 vs.\ .49); domain-specific subword embeddings beat Wikipedia embeddings by 14\,\% at $k=1$.

\cf{User study} Between-subjects, four surveys (Retrieval, SemVec, PriBot, Random), 700 MTurk workers (success rate $>$95\,\%, \$2 for about 14 minutes), each rating 15 QA pairs plus two attention checks, minimum 15 seconds per page, a Cloze reading test to remove low-skill responses; five-point relevance scale, an answer counts as relevant if at least 50\,\% of at least seven raters say so \pc{harkous2018}{14}{We used a between-subject design by constructing four surveys, each corresponding to a different evaluation condition}. Result: top-1/2/3 relevance for PriBot .70/.78/.89 vs.\ Retrieval .46/.71/.79, SemVec .48/.71/.85, Random .37/.59/.76; the authors note that predictive accuracy and perceived usefulness differ, as in recommender systems \pc{harkous2018}{15}{The respondents regarded at least one of the top-3 answers as relevant for 89\% of the questions, with the first answer being relevant in 70\% of the cases}.

\cf{Limits (authors and card)} Bound to the OPP-115 taxonomy \pc{harkous2018}{15}{Polisis might be limited by the employed privacy taxonomy}; misclassification should be reported with confidence scores; not robust to adversarially worded policies \pc{harkous2018}{15}{Polisis is not guaranteed to be robust in handling an adversarially constructed privacy policy}; not a replacement for the legally binding text \pc{harkous2018}{15}{We also want to stress the fact that Polisis is not intended to replace the legally-binding privacy policy}. Card additions: the headline icon accuracy compares the system with the \emph{union} of three annotators' labels on only 50 test policies (accuracy depends on a lenient reference); the icon rules are the authors' interpretation; the weakest classes (retention, choice) are those most relevant to consent quality; the Twitter test set is small (120) and annotated by two authors who are also developers; MTurk raters judge topical relevance, not legal correctness; policies are English, web-based and from before GDPR; the neural models predate LLMs, so absolute numbers cannot be transferred.

\cf{Use for this thesis} (1) The evaluation recipe for stage L: hold out policies, compare machine labels with the aggregated human labels on the same unit, report accuracy \emph{and} a chance-corrected $\kappa$ per task, state the aggregation rule for the human reference, and compare to a baseline. (2) A reminder that high accuracy against a union or majority label can mask weakness on rare or difficult categories, which the thesis should report per code. (3) Evidence on where classifiers fail on policies (retention period, user choice), which suggests which codes will need extra human checking in an LLM pipeline. (4) A methodology for reporting both human--human and human--machine agreement on identical categories (the tweet and ground-truth $\kappa$ values), and the point that predictive accuracy and perceived usefulness are not the same. (5) It supports the design choice of treating automatic annotation as decision support with a human in the loop. Treat the OPP-115 pipeline as a baseline to cite, not as something to reproduce.
\end{papercard}
